%%%% ABSTRACT
%%
%% Versão do resumo para idioma de divulgação internacional.

\begin{abstractutfpr}%% Ambiente abstractutfpr
This work investigates the feasibility of distributing a proof-of-concept application based on the Notification-Oriented Paradigm (NOP) in a Kubernetes cluster, considering aspects such as component decoupling, scalability, resource consumption, and execution organization in a distributed environment. The study is grounded in the context of computationally demanding systems and in the need to assess whether NOP principles, especially the reduction of structural and temporal redundancies, can support a more suitable architecture for distribution. To this end, the theoretical basis relies on NOP Query, a notification-based method applied to continuous-stream query processing, whose proposal shows that small and collaborative entities can operate with low latency and tractable computational cost in data-intensive scenarios. Based on this foundation, the investigation analyzes how the components of a proof-of-concept application can be containerized and orchestrated by Kubernetes, enabling independent deployment, communication among services, and horizontal scaling. The expected findings indicate that combining NOP and Kubernetes is technically plausible for this type of application, although limitations still exist regarding state management, inter-component communication, and memory consumption in notification-based structures. It is concluded that distribution in Kubernetes is compatible with the NOP proposal and may expand its applicability in modern systems, provided that the architecture is designed with attention to coordination costs and to the constraints of the proof-of-concept scope.
\end{abstractutfpr}
